% Part: computability
% Chapter: recursive-functions
% Section: notations-pr-functions

\documentclass[../../../include/open-logic-section]{subfiles}

\begin{document}

\olfileid{cmp}{rec}{not}
\olsection{ترميزات العودية البدائية}

إذا ضبطنا الدوال العودية البدائية بتعريف استقرائي دقيق، أمكننا أن نسلك
في وصفها سبيلًا منتظمًا؛ ومن ذلك أن نجعل لكل دالة «ترميزًا» على
الوجه الآتي. نضع الرموز $\Zero$ و$\Succ$ و$\Proj{n}{i}$ للصفر
واللاحق والإسقاطات. فإذا كانت $h$ مركبة من دالة~$f$ ذات
$k$ حجج ومن دوال $g_0$، …،~$g_{k-1}$ ذات $n$ حجة، وكنا قد جعلنا
الترميزات $F$ و$G_0$، …،~$G_{k-1}$ لهذه الدوال، جاز لنا
أن نستعمل رمزًا جديدًا، هو $\fn{Comp}_{k,n}$، فنرمز إلى الدالة $h$ بـ
$\fn{Comp}_{k,n}[F,G_0,\dots,G_{k-1}]$.

وكذلك نصنع في الدوال المعرفة بالعودية البدائية. فإذا
كانت الدالة~$h$ ذات $k+1$ حجة معرفة بالعودية البدائية من الدالة~$f$ ذات
$k$ حجج والدالة~$g$ ذات $k+2$ حجة، وكان الترميزان المسندان إلى $f$
و~$g$ هما $F$ و~$G$ على الترتيب، فالترميز الذي نجعله لـ~$h$ هو
$\fn{Rec}_k[F,G]$.
% تصحيح صياغة كلاسيكية (C214-01): رُبط الشرط بجميع مقدماته ثم بجوابه.

وقد تقدم أن تعريف دالة الجمع بالعودية البدائية هو:
\begin{align*}
  \Add(x_0, 0) & = \Proj{1}{0}(x_0) = x_0\\
  \Add(x_0, y+1) & = \Succ(\Proj{3}{2}(x_0, y, \Add(x_0, y))) = \Add(x_0, y) +1
\end{align*}
فـ$\Proj{1}{0}$ هنا في موضع~$f$، و
$\Succ(\Proj{3}{2}(x_0,y,z))$ في موضع~$g$؛ وترميزها
$\fn{Comp}_{1,3}[\Succ,\Proj{3}{2}]$ لأنها ناتجة من تعريف دالة
بالتركيب من الدالة~$\Succ$ الأحادية الحجة والدالة~$\Proj{3}{2}$
الثلاثية الحجج. وعلى هذا يكون ترميز دالة الجمع
\[
\fn{Rec}_1[\Proj{1}{0},\fn{Comp}_{1,3}[\Succ,\Proj{3}{2}]].
\]
ولهذه الترميزات فائدة في مواضع، منها تعداد الدوال العودية البدائية.

\begin{prob}
  اكتب ترميز~$\Mult$ العودي البدائي تامًا.
\end{prob}

\end{document}
